\documentclass[12pt]{article}

\usepackage{graphicx}
\usepackage{fancyhdr}
\usepackage{geometry}
\usepackage{xcolor}
\usepackage{enumitem}
\usepackage{url}

\geometry{margin=1in}

\title{CSE 123A - Demo Script\\Water Level Monitoring System}
\author{Group 8}
\date{}

\begin{document}
\pagestyle{fancy}
\maketitle
\fancyfoot{}
\fancyhead{}
% \fancyfoot[R]{\thepage}

%% --------------------------------------------------------
\section*{Part 1: Introduction}
\noindent\textbf{Presenter:} TODO\\  
\noindent\textbf{Estimated Time:} 30 seconds

\noindent\textbf{Action:} Hold up or gesture toward the water pitcher.

\bigskip
\noindent\textit{Script:}

\begin{quote}
``Has this ever happened to you? You come home after a long day, you're thirsty, and you go to grab the water filter pitcher---only to find it completely empty. Nobody refilled it.
This is a common problem in shared living spaces---dorms, family homes, offices. Everyone uses the filtered water, but nobody knows when it's running low until it's too late.
Our project solves this. We built a smart water level monitoring system that sits underneath your water filter pitcher, continuously tracking how much water is left, and sends a notification to everyone in the household when it's time to refill.''
\end{quote}

%% --------------------------------------------------------
\section*{Part 2: Overview }
\noindent\textbf{Presenter:} Reid\\  
\noindent\textbf{Estimated Time:} 1 minute

\noindent\textbf{Action:} Point at the various system components as they are mentioned.

\bigskip
\noindent\textit{Script:}

\begin{quote}
``Our design has three main components: the smart base, the web server, and the mobile app.

The smart base sits beneath the water pitcher and measures the weight of the pitcher to determine how much water is left. This data is sent to our web server over WiFi using HTTP packets, where it's stored in a database and processed. 

The mobile app is connected to the web server and receives push notifications when the water level drops below a certain threshold. The app also allows users to view the current water level and configure various settings.

\textit{[Switch to prototype diagram.]}
Today, we will be demonstrating the prototype of our design. We are simulating the connection between the smart base and the web server through two laptops: one connected to the microcontroller to read the sensor data, and another sending simulated weight data to the web server. Our mobile app is currently in the form of a webapp on a laptop, but the functionality is the same as a mobile app would be.'' 


''
\end{quote}

%% --------------------------------------------------------
\section*{Part 3: The Hardware }
\noindent\textbf{Presenter:} TODO\\  
\noindent\textbf{Estimated Time:} TODO

\noindent\textbf{Action:} Pick up the prototype base and point out each component as mentioning it.

\bigskip
\noindent\textit{Script:}

\begin{quote}
``Let's start with the hardware. The system has three main components:

\textbf{First, the load cell.} This is a strain-gauge weight sensor. When weight is applied, it bends slightly, and that bending changes its electrical resistance. The change is very tiny---so we use the \textbf{HX711 amplifier board} to boost the signal into something our microcontroller can read. The HX711 gives us 24-bit resolution, which means we can detect very small changes in weight.

\textbf{Second, the microcontroller.} This is the brain of our system. It reads the weight data from the HX711, converts it into a water level percentage, and sends that data over WiFi to our cloud server.

\textbf{Third, the mounting base.} We designed and 3D-printed two circular plates, that sandwich the load cell. The pitcher sits on top. We designed these with M4 screw holes so everything stays secure and aligned.''
\end{quote}


%% --------------------------------------------------------
\section*{Part 4: Live Demo -- Weighing }
\noindent\textbf{Presenter:} TODO\\  
\noindent\textbf{Estimated Time:} TODO

\noindent\textbf{Action:} Place the pitcher on the platform. Have the a computer open showing the readings of the sensor.

\bigskip
\noindent\textit{Script:}

\begin{quote}
``Now let's see it in action.
% maybe add calibration before?
I'm going to place this water pitcher on the sensor base. The load cell is now measuring the total weight. Our microcontroller reads that value, subtracts the known weight of the empty pitcher, and calculates the percentage of water remaining.

\textit{[Point to the web dashboard.]}
% ( with fake data)
Over here on our web dashboard, you can see the water level updating. Right now it's showing [read current percentage]. The pitcher graphic fills up to match the actual level, and you can see the exact percentage and the last time the reading was updated.

Let me pour some water out to simulate usage.

\textit{[Pour water out of the pitcher into a cup.]}
% fake data 
Watch the dashboard---you can see the percentage dropping as I pour. The system polls the database every five seconds, so there's a short delay, but it tracks the change.''
\end{quote}


% show how software works

\section*{Part 5: Software Demo}
\noindent\textbf{Presenter:} TODO\\  
\noindent\textbf{Estimated Time:} TODO

\noindent\textbf{Action:} Open laptop to the web server to demonstrate navigating the page

\bigskip
\noindent\textit{Script:}

\begin{quote}
``The initial prototype uses a web dashboard to display important device information. This dashboard is also where the initial configuration of the product happens. When a user first starts up this device, they need to calibrate it with empty and full weights. 

\textit{[Highlight calibration buttons]} Both an empty and a full value are required for the product to work as intended. In the startup sequence the user will be prompted for an empty pitcher of water to be placed on the base. After pressing the calibration button they will be asked to fill it up and return it to the base to calibrate the full value. 

Full and empty values are adjustable and configurable at any time, so if a user buys a new filter that holds a larger quantity of water, they can simply change the calibration values and continue using the device as intended. 

When the values are calibrated, the system is able to use the data sent by the microcontroller to adjust the visible water level."
\end{quote}

% show how notifications work
\section*{Part 6: Notification Demo}
\noindent\textbf{Presenter:} TODO\\  
\noindent\textbf{Estimated Time:} TODO

\noindent\textbf{Action:} Keep the dashboard visible and reset the water level to above the low-water threshold 20\%.

\bigskip
\noindent\textit{Script:}

\begin{quote}
``Now we will demonstrate the notification feature and how we can subscribe to alerts for when the water level is running low.

\textit{[Point to the current percentage on the dashboard.]}
Right now, our low-water threshold is set to 20\% and the pitcher level is above the low-water threshold, so no alert is sent. Once the level drops below that value, the app will trigger a push notification to subscribed devices. But first, we have to subscribe to notifications. 

\textit{[Point to "Enable Notifications" button on the dashboard.]} This button will prompt you to allow notifications for our web app. Once you allow that, your device will be subscribed and ready to receive alerts.

\textit{[Click on "Enable Notifications" button and allow]} Now, we will measure the pitcher when it is near empty. 

\textit{[Simulate water level to less than 20\%]} You can see the dashboard now reads [read value below 20\%]. Within a few seconds, the notification should appear on the top right corner of your screen.

The message says that the filter is running low and needs to be refilled. After the pitcher is refilled and the percentage rises, the system returns to normal and waits for the next low-water event.''
\end{quote}

% Q&A

\section*{Part 7: Questions}
\noindent\textbf{Presenter:} TODO (prof suggested same as intro person)\\  
\noindent\textbf{Estimated Time:} 5 minutes

\noindent\textbf{Action:} Leave all components visible on the center of the demo table, ready to be moved or pointed at for any questions regarding hardware. 


% conclusion
\end{document}

